\section{Ultra-high-energy cosmic rays as a probe of a cubic spacetime lattice}
\label{sec:uhecr}

This section covers the observational side of the proposal by Beane, Davoudi
and Savage~\cite{Beane2014constraints}. It reviews what measurements of
ultra-high-energy cosmic rays (UHECRs) say about their lattice scenario,
separates what current data constrain from what they leave open, and ends with
a specific analysis that could be run on public data. The lattice theory and
the Lorentz-violation bounds are treated in Sec.~\ref{sec:lattice}.

\subsection{The prediction and its premises}
\label{sec:uhecr-prediction}

Beane \emph{et al.} assume a simulation on a hyper-cubic space-time grid with
spacing $b$, starting from an unimproved Wilson action~\cite{Beane2014constraints}.
The lattice dispersion relations cap the energy at $E_{\max}\sim 1/b$ for
fermions and bosons. Setting this cap equal to the Greisen--Zatsepin--Kuzmin
(GZK) cutoff~\cite{Greisen1966end,Zatsepin1966upper} gives
$b\sim10^{-12}$\,fm, that is $b^{-1}\sim10^{11}$\,GeV. The authors turn this into
the bound $b^{-1}\gtrsim10^{11}$\,GeV, their most stringent
one~\cite{Beane2014constraints}. They also argue that if the lattice itself
supplies the cutoff, ``the angular distribution of the highest energy
components would exhibit cubic symmetry in the rest frame of the lattice'',
and that for smaller $b$ the cubic signal would fade while the GZK mechanism
takes over~\cite{Beane2014constraints}. Their Eq.~(18) supports this at leading
order. There the lattice shifts the $\Delta$-resonance threshold of
proton--CMB-photon collisions in proportion to the combination
$Y_4^0+\sqrt{5/14}\,(Y_4^4+Y_4^{-4})$, taking the lattice rest frame to be the
CMB frame~\cite{Beane2014constraints}.

The prediction has three limits that matter here.
(i)~It is qualitative. The paper gives neither the size of the flux anisotropy
as a function of $b$ nor a sky template.
(ii)~It assumes that the observed steepening is the GZK cutoff, and it treats
the primaries through single-particle dispersion relations. The paper does not
discuss heavy nuclei as primaries~\cite{Beane2014constraints}. We add that a
nucleus of mass number $A$ carries only about $E/A$ per nucleon, so identifying
$E_{\max}$ with the observed endpoint depends on composition.
(iii)~The arXiv version puts the GZK scale at ``$\sim6\times10^{20}$\,eV''
\cite{Beane2014constraints}. That is ten times the threshold of
$\sim6\times10^{19}$\,eV quoted by the experiments~\cite{HiRes2008first}, and
we do not reproduce it.

We checked numerically which spherical-harmonic degrees admit functions
invariant under the full cubic group $O_h$ (48 elements). Up to $\ell=12$ they
are $\ell=0,4,6,8,10$ and $12$ (two at $\ell=12$). The combination in Beane
\emph{et al.}'s Eq.~(18) is the $\ell=4$ invariant, and
$Y_6^0-\sqrt{7/2}\,(Y_6^4+Y_6^{-4})$ is the $\ell=6$ one. It follows that a
lattice-aligned pattern has no dipole and no quadrupole: its leading content
is at $\ell=4$ and $\ell=6$.

\subsection{The flux suppression and the composition question}
\label{sec:uhecr-spectrum}

HiRes reported a suppression at $\log_{10}(E/\mathrm{eV})=19.75\pm0.04$. It
observed 13 events above $10^{19.8}$\,eV where an unbroken spectrum predicted
43.2, a Poisson deficit equivalent to 5.3 standard deviations
\cite{HiRes2008first}. The Pierre Auger Observatory rejected a constant slope
above $4\times10^{19}$\,eV at 6 standard deviations, while cautioning that
identifying the cause would require knowledge of the mass
composition~\cite{PierreAuger2008observation}. Auger's 2020 spectrum uses
215{,}030 events and an exposure of $(60{,}400\pm1{,}810)$\,km$^2$\,sr\,yr
\cite{PierreAuger2020measurement}. It places the suppression at
$E_{34}=(46\pm3\pm6)\times10^{18}$\,eV with spectral index
$\gamma_4=5.1\pm0.3\pm0.1$, and finds a new softening (the ``instep'') at
$E_{23}=(13\pm1\pm2)\times10^{18}$\,eV~\cite{PierreAuger2020features}. The
measured $E_{1/2}=(2.2\pm0.1\pm0.3)\times10^{19}$\,eV differs from the
$5.3\times10^{19}$\,eV expected for uniformly distributed, proton-only sources,
a scenario Auger describes as disfavored~\cite{PierreAuger2020features}.

Mass composition decides how the endpoint should be read. In Auger's benchmark
mixed-composition model, the steepening above $\approx5\times10^{19}$\,eV comes
from the maximum acceleration energy of the heaviest nuclei combined with the
GZK effect. Auger adds that a subdominant light component is not excluded and
that other scenarios remain viable~\cite{PierreAuger2020features}. A combined
fit of spectrum and composition prefers sources with a mixed composition, a
hard spectrum and a low rigidity cutoff~\cite{PierreAuger2023constraining}.
Auger has also inferred $X_{\max}$ event by event from surface-detector data up
to 100\,EeV~\cite{PierreAuger2025inference}. That analysis finds a composition
that becomes ``increasingly heavier and purer'' and is incompatible with a large
light fraction between 50 and 100\,EeV. It concludes that the suppression
``cannot be entirely ascribed to effects of extragalactic propagation''
\cite{PierreAuger2025inference}. Telescope Array (TA) hybrid data are
compatible with QGSJet\,II-04 protons at 95\% CL in all energy bins, after
systematic shifts. At the highest energies, where statistics are sparse, other
elements are also compatible~\cite{TelescopeArray2018depth}. In short, the
premise that the observed endpoint is the GZK cutoff of a proton flux is
disputed by current data. This weakens the inference from ``the spectrum ends
near $10^{20}$\,eV'' to a statement about $b$, whether the end is attributed to
the GZK effect or to a lattice.

\subsection{Arrival directions at each angular scale}
\label{sec:uhecr-anisotropy}

\paragraph{Dipole.}
Auger observed a dipole above 8\,EeV at more than $5.2\sigma$ (post-trial),
using $3\times10^4$ events from an exposure of 76{,}800\,km$^2$\,sr\,yr. Its
amplitude is $6.5^{+1.3}_{-0.9}\%$ towards $(\alpha,\delta)=(100\pm10^\circ,
-24^{+12}_{-13}{}^\circ)$~\cite{PierreAuger2017observation}. With 19 years of data
(49{,}678 events $\geq8$\,EeV), the right-ascension modulation reaches $6.8\sigma$.
The 3D dipole is $7.4^{+1.0}_{-0.8}\%$ towards
$(97\pm8^\circ,-38\pm9^\circ)$, and it does not vary in time: the rate of change
is below 0.3\% per year at 95\% CL~\cite{PierreAuger2024large}. In the joint
dipole-plus-quadrupole fit, the quadrupole amplitude for $E\geq32$\,EeV is
$Q=0.13\pm0.06$, and the quadrupolar components are not
significant~\cite{PierreAuger2024large}.

\paragraph{Intermediate scales.}
A starburst-galaxy sky model fits 5{,}514 Auger events above 20\,EeV better than
isotropy at $4.0\sigma$~\cite{PierreAuger2018indication}. With the full Phase~1
set of 2{,}635 events above 32\,EeV (exposure 122{,}000\,km$^2$\,sr\,yr), Auger
finds evidence at the $4\sigma$ level for an excess on a $\sim15^\circ$
(Gaussian) to $\sim25^\circ$ (top-hat) scale above $\sim40$\,EeV. The most
significant blind excess lies in the Centaurus region, with $5.4\sigma$ local
significance and a 3\% post-trial probability~\cite{PierreAuger2022arrival}. In
the north, TA reported a hotspot above 57\,EeV with local significance
$5.1\sigma$ and chance probability $3.7\times10^{-4}$ ($3.4\sigma$)
\cite{TelescopeArray2014indications}. The joint Auger--TA update finds the
starburst correlation to be the most significant, at $4.2\sigma$ post-trial
\cite{PierreAugerTA2025new}. TA also recorded a single event of
$244\pm29\,(\mathrm{stat})\,^{+51}_{-76}\,(\mathrm{syst})$\,EeV whose direction
points back to a void~\cite{TelescopeArray2023extremely}.

\paragraph{Higher multipoles, including $\ell=4$ and $\ell=6$.}
Table~\ref{tab:uhecr-multipoles} lists the analyses that measured multipoles
beyond the quadrupole. Auger observes only part of the sky, so it cannot
estimate individual $a_{\ell m}$ with useful resolution once $\ell_{\max}>2$.
It reconstructs $C_\ell$ instead, assuming the anisotropies are one
realization of a statistically isotropic random
process~\cite{PierreAuger2024large,PierreAuger2017multiresolution}. The 2024
spectrum (with $\ell=1$--20 in each of four energy bins) shows no significant
$\hat C_\ell$ apart from the dipole. Two excursions, $\hat C_{17}$ at 4--8\,EeV
and $\hat C_8$ at 16--32\,EeV, have penalized chance probabilities of 3.3\% and
26.5\%~\cite{PierreAuger2024large}. The 99\% CL upper limits on $C_\ell$ are
shown only graphically; no numerical $C_4$ or $C_6$ limits are
tabulated~\cite{PierreAuger2024large}. With full-sky Auger+TA data, the
$a_{\ell m}$ have been estimated up to $\ell=20$ above $10^{19}$\,eV, with no
evidence for anisotropy~\cite{PierreAugerTA2014searches}. A newer harmonic-space
analysis over $\ell\le20$ with a scan in energy threshold finds the quadrupole
to be the most significant multipole in every case. Its largest
auto-correlation excursion is at $\ell=2$ above 41\,EeV (Auger energy scale),
at $4.2\sigma$ pre-trial and $2.1\sigma$ post-trial~\cite{PierreAugerTA2025new}.
An independent likelihood reanalysis of the public 2017 Auger events likewise
finds nothing beyond the dipole~\cite{Ahlers2018searching}.

\begin{table}[t]
\centering
\small
\caption{Published analyses that measured UHECR multipoles beyond $\ell=2$,
all of which therefore contain $\ell=4$ and $\ell=6$. None forms an
$O_h$-invariant combination or scans lattice orientations.}
\label{tab:uhecr-multipoles}
\begin{tabular}{p{0.24\linewidth}p{0.2\linewidth}p{0.17\linewidth}p{0.29\linewidth}}
\hline
Analysis & Data, energy & $\ell$ range & Result beyond the dipole\\
\hline
Auger 2017~\cite{PierreAuger2017multiresolution} & Auger, 4--8 and $\geq8$\,EeV & $C_\ell$, $\ell\le64$ & no deviation from isotropy\\
Auger 2024~\cite{PierreAuger2024large} & Auger, 4 bins $\geq4$\,EeV & $C_\ell$, $\ell\le20$ & $\hat C_{17}$, $\hat C_8$ not significant after penalty; 99\% UL graphical\\
Auger+TA 2014~\cite{PierreAugerTA2014searches} & full sky, $>10$\,EeV & $a_{\ell m}$, $C_\ell$, $\ell\le20$ & no significant deviation\\
Auger+TA 2025~\cite{PierreAugerTA2025new} & full sky, threshold scan & $C_\ell$ (auto, cross), $\ell\le20$ & $\ell=2$ most significant (2.1$\sigma$ post-trial)\\
Ahlers 2018~\cite{Ahlers2018searching} & public Auger, $>8$\,EeV & pseudo-$C_\ell$ & consistent with isotropy\\
\hline
\end{tabular}
\end{table}

\subsection{Magnetic smearing}
\label{sec:uhecr-gmf}

Deflections scale with charge and with the integrated transverse field, and
inversely with energy~\cite{PierreAuger2017observation}. At
$E\approx10$\,EeV, estimates of the mean charge range from $Z=1.7$ to $5$, and
typical Galactic deflections are a few tens of degrees at $E/Z=10$\,EeV
\cite{PierreAuger2017observation}. In the Jansson--Farrar (JF12) model, a
60\,EeV proton is deflected by $5.2^\circ$ on average. A quarter of the sky has
deflections below $2.2^\circ$, and deflections in the south are about 60\%
larger than in the north~\cite{Jansson2012new}. The UF23 ensemble of eight
coherent-field models~\cite{Unger2024coherent} requires a rigidity of
$R\geq20$\,EV for deflections below $20^\circ$ over half the sky in all models.
If directions are corrected for the modelled field, the requirement falls to
11\,EV; the authors call this estimate indicative because random fields are
not included~\cite{Unger2024coherent}. The model-to-model spread is smaller than
the deflections themselves over most of the sky~\cite{Unger2024coherent}.

For iron-like charges near the suppression energy, these rigidities imply
deflections of tens of degrees. That is comparable to the $\sim45^\circ$ angular
scale of an $\ell=4$ pattern, so the Galactic field would distort a
lattice-aligned pattern in a direction-dependent way. A search must therefore
either forward-model the distortion or work at the highest rigidities, where
the pattern survives.

\subsection{What the existing data do and do not constrain}
\label{sec:uhecr-status}

$C_\ell$ does not change under rotations. A cubic pattern of \emph{any}
orientation therefore adds power to $C_4$ and $C_6$, and the non-detections in
Table~\ref{tab:uhecr-multipoles} already bound the size of such a pattern
without any orientation scan. The bound is only implicit, for four reasons.
(a)~No paper quotes numerical $C_4$ or $C_6$ limits above the suppression
energy.
(b)~The partial-sky $C_\ell$ reconstruction assumes statistical isotropy, which
a fixed lattice-aligned sky violates~\cite{PierreAuger2017multiresolution}.
(c)~$C_4$ and $C_6$ each spread their statistical weight over $2\ell+1=9$ and
13 directions and ignore the fixed phase relation that a cubic pattern imposes
between $\ell=4$ and $\ell=6$. Blind scans over many multipoles lose
sensitivity: for a 5\% dipole in 14{,}000 simulated events, a blind
$\ell\le64$ power-spectrum search detects it at 99\% CL in only 25\% of skies
at $\delta=-30^\circ$~\cite{PierreAuger2017multiresolution}.
(d)~Turning an amplitude limit into a limit on $b$ requires the anisotropy
amplitude as a function of $b$ and rigidity, which Beane \emph{et al.} do not
provide~\cite{Beane2014constraints}.

\paragraph{Novelty check.}
We searched for any analysis that tests UHECR (or other astroparticle) arrival
directions for a lattice-aligned, cubic or $O_h$-symmetric anisotropy, or that
confronts the ``numerical simulation'' hypothesis with such data. We checked
every work that cites Ref.~\cite{Beane2014constraints} in INSPIRE (13 records)
and Semantic Scholar (61 records), ran full-text queries in INSPIRE, and ran
arXiv and general web searches. The full log is in the project repository. We
found no such analysis. Citing works are theory, philosophy or popular
accounts. Among them, Ref.~\cite{Vazza2025astrophysical} uses the UHECR energy
only as a resolution requirement for a simulator, and
Ref.~\cite{Mlodinow2025bounds} bounds a cubic-lattice quantum cellular
automaton with photon timing and laboratory data. The closest data analyses test
\emph{direction-dependent Lorentz violation}, not a cubic flux pattern:
vacuum-Cherenkov bounds on anisotropic modified-Maxwell coefficients from the
directions of UHECR events~\cite{Klinkhamer2008ultrahigh}, and a sidereal-modulation
search in IceCube atmospheric neutrinos~\cite{IceCube2010search}. Auger's own
Lorentz-violation study uses direction-independent dispersion
relations~\cite{PierreAuger2022testing}. To our knowledge, no search for an
$O_h$-symmetric UHECR anisotropy has been published. NASA ADS and Google Scholar
could not be queried automatically, so this statement carries that caveat.

\subsection{Public data}
\label{sec:uhecr-data}

Auger opened its Open Data Portal in February 2021 with 10\% of Phase~I data.
The release now holds 81{,}121 showers, including 25{,}086 events from the
1500\,m array: vertical above 2.5\,EeV and inclined ($60^\circ$--$80^\circ$)
above 4\,EeV. The data come as JSON files per event and a CSV
summary~\cite{PierreAuger2025pierre,PierreAuger2024pierre}. In June 2023 the
collaboration approved raising the public fraction to 30\% of main-array events
above 2.5\,EeV from 2004--2022 (about 24{,}000\,km$^2$\,sr\,yr), with release
planned for late 2025~\cite{PierreAuger2025expanding}. When we checked the
portal and Zenodo on 27 September 2026, the current version was still the 10\%
release of March 2024~\cite{PierreAuger2024pierre}. Our count of the public
summary files gives only 86 events above 40\,EeV. Much more useful for this
problem is the \emph{complete} Phase~1 list of 2{,}635 events above 32\,EeV. It
was published with Ref.~\cite{PierreAuger2022arrival} and gives, for each event,
its time, local and equatorial coordinates, energy and cumulative
exposure~\cite{PierreAuger2022material}. Our count of that file gives 1{,}387
events at $E\geq40$\,EeV and 411 at $E\geq57$\,EeV. Northern coverage is thinner
in public form: TA published the 72 events above 57\,EeV from its first five
years~\cite{TelescopeArray2014indications}, and Auger catalogued its 100
highest-energy events (78--166\,EeV)~\cite{PierreAuger2023catalog}. An external
all-scale reanalysis of public Auger events has already been
published~\cite{Ahlers2018searching}, so such work is practical.

\subsection{A proposed search (not a result)}
\label{sec:uhecr-proposal}

The following analysis is feasible with the public data above and addresses
points (b) and (c) of Sec.~\ref{sec:uhecr-status}. We present it as a proposal;
it has not been carried out.

\emph{Model.} The relative flux is written as
$\Phi(\hat n)\propto\omega(\hat n)\,[1+\vec d\cdot\hat n+\epsilon\,
K(R^{-1}\hat n)+\dots]$. Here $\omega$ is the directional exposure, $\vec d$ is a
nuisance dipole (plus quadrupole) with priors from
Refs.~\cite{PierreAuger2024large,PierreAuger2017observation}, and
$K=\cos\chi\,\hat K_4+\sin\chi\,\hat K_6$ is a unit-normalized combination of
the $\ell=4$ and $\ell=6$ $O_h$ invariants. The rotation $R\in SO(3)/O$ fixes
the lattice orientation; the quotient is by the 24 proper rotations of the
cube, since even-$\ell$ functions are automatically inversion-symmetric.

\emph{Statistic.} The test statistic is the likelihood ratio
$T=\max_{R,\chi}\,2\ln[\mathcal L(\epsilon,R,\chi,\vec d)/\mathcal L(0,\vec d)]$
from an unbinned likelihood over events. For a template-free cross-check, the
``cubicness'' of the measured $\ell=4$ multipole, $\max_R|\langle\hat K_4,
D^{(4)}(R)\,a_4\rangle|^2/\sum_m|a_{4m}|^2$, equals 1 for an exactly cubic
$\ell=4$ component and has a computable null distribution. The two degrees are
tied together by sharing a single $R$, which is what distinguishes the test
from separate $C_4$ and $C_6$ measurements.

\emph{Look-elsewhere correction.} Isotropic-plus-dipole mock skies drawn with
the actual exposure give the null distribution of $T$. This distribution already
includes the three-parameter orientation scan and the $\chi$ scan. Energy
thresholds (for example 32, 40 and 57\,EeV on the Auger scale) should be fixed
in advance, or else penalized in the same simulations.

\emph{Exposure and coverage.} With Auger data alone, the unobserved northern
cap correlates the $a_{\ell m}$. The likelihood in pixel or event space handles
this directly and avoids the statistical-isotropy assumption behind pseudo-$C_\ell$
deconvolution. Adding the TA list~\cite{TelescopeArray2014indications} needs a
cross-calibration of energy scales of the kind used in the joint
analyses~\cite{PierreAugerTA2014searches,PierreAugerTA2025new}.

\emph{Magnetic forward model.} Lattice templates defined outside the Galaxy
should be propagated through each of the eight UF23 variants~\cite{Unger2024coherent}
and JF12~\cite{Jansson2012new} for a grid of rigidities. The spread of results
is a systematic uncertainty, and the unperturbed template is the
high-rigidity limit.

\emph{Validation.} Before touching the data, one should inject cubic patterns
of known $(\epsilon,R,\chi)$ into mock skies that also contain the measured
dipole and a Centaurus-like hotspot~\cite{PierreAuger2022arrival}. The checks
are that the orientation is recovered, that the false-positive rate is correct,
and what the detection power is as a function of $\epsilon$. An idealized
estimate that we verified by Monte Carlo sets the scale: for $N$ events
distributed isotropically with uniform full-sky exposure, the amplitude
estimator of one unit-normalized harmonic has standard deviation
$\sqrt{4\pi/N}$, about 0.095 for $N=1{,}387$. Partial coverage, the nuisance
dipole and the orientation scan will increase the threshold for detection.

\emph{Interpretation.} A null result would bound $\epsilon$ as a function of
energy threshold and assumed rigidity. Relating $\epsilon$ to $b$ needs a
calculation of the anisotropy amplitude that the literature does not yet
contain (Sec.~\ref{sec:lattice}). The bound also does not apply to
discretizations with other symmetries or with improved actions, which Beane
\emph{et al.} themselves consider likely~\cite{Beane2014constraints}.
